% ===========================================================================
\section{Emission Factors: Table 2.5, Not Table 2.2}
\label{sec:factors}
% ===========================================================================

\subsection{Which of the four tables}

IPCC Volume~2 Chapter~2 publishes four parallel default tables for stationary combustion---
energy industries, manufacturing, commercial, and residential---and this sector must use
Table~2.5, Residential and Agriculture/Forestry/Fishing.

\begin{table}[htbp]
\centering
\footnotesize
\begin{tabular}{@{}lrrr@{}}
\toprule
\textbf{Natural gas} & \textbf{\COtwo{}} & \textbf{\CHfour{}} & \textbf{\NtwoO{}} \\
 & \multicolumn{3}{c}{kg/TJ} \\
\midrule
Table 2.5, residential        & 56\,100 & 5 & 0.1 \\
Table 2.2, energy industries  & 56\,100 & 1 & 0.1 \\
\midrule
\textbf{Kerosene} & & & \\
Table 2.5, residential        & 71\,900 & 10 & 0.6 \\
Table 2.2, energy industries  & 71\,900 &  3 & 0.6 \\
\bottomrule
\end{tabular}
\caption{\COtwo{} is identical because it depends only on the carbon in the fuel. \CHfour{} is
not: residential natural gas is five times the utility value and residential kerosene more than
three times, because a small appliance burns less completely than a power-station boiler.}
\label{tab:tables}
\end{table}

Taking Table~2.2 here would understate household methane by \SI{80}{\percent} while leaving
\COtwo{} exactly correct---and \COtwo{} is \SI{76}{\percent} of the sector's reported
emissions, so every obvious check would still pass. This is one of four instances of the same
trap documented across this inventory \citep{ahmad2026national}; the power sector's module
carries the mirror-image warning, and a test pins the direction in both.

\subsection{Biomass factors, and why their bands are wide}

The dung and firewood streams carry uncertainties above \SI{150}{\percent} on their methane and
nitrous oxide. That is not conservatism: cookstove emission factors genuinely span that range,
because combustion completeness depends on stove design, fuel moisture, feeding rate and
operator behaviour, and measurement campaigns disagree substantially.

The consequence is that the biomass streams that supply most of the sector's energy are also its most
uncertain emission streams, which is why the sector's primary band---the grouped Approach-1
figure of $\pm\SI{23.1}{\percent}$ (Section~\ref{sec:results})---is wide despite three of
the nine streams being metered: the metered gas, LPG and kerosene core is
power-station-grade, and everything around it is not.

\subsection{Fossil and biogenic methane in one category}

This sector is the only one in the inventory that emits both. Dung and firewood methane is
biogenic and takes the AR5 GWP of \num{28}; natural gas, LPG and kerosene methane is fossil and
takes \num{30}. Every methane-emitting component declares which it is, and a contract test
asserts both that the declaration is present and that the \COe{} conversion used the
corresponding value.

The difference is small in this sector---methane totals \SI{5.18}{\mega\tonne} \COe{} of which
\SI{0.05}{} is fossil---but the same field distinguishes landfill methane from combustion
methane elsewhere in the inventory, and a category emitting both is exactly where a missing
flag would be least visible.

\subsection{No separate oxidation factor}

The 2006 Guidelines fold complete oxidation into the emission factor; Table~1.4 carries an
explicit oxidation column of \num{1.0} for every fuel. Applying the 1996-vintage \num{0.98} on
top would double-discount. The waste sector in this inventory \emph{does} carry an oxidation
factor, because open burning genuinely is incomplete, and tests pin both conventions.
